\SourceLeaf{243}%
biegsamen, undehnbaren Körper auf, von dessen Dimensionen eine als verschwindend gilt, so hängen die Eigenschaften der Fläche zum Teil von der Form ab, in die man sie sich gebracht denkt. Zum Teil sind sie absolut und bleiben unverändert, in welche Form man die Fläche auch biegt. Die Untersuchung dieser letzteren Eigenschaften eröffnet der Geometrie ein neues und fruchtbares Feld. Zu ihnen gehören das Krümmungsmaß und die Integralkrümmung in dem Sinne, in dem wir diese Ausdrücke gebrauchen. Hierher gehören ferner die Theorie der kürzesten Linien und mehrere andere Gegenstände, deren Behandlung wir uns für später vorbehalten. Bei dieser Betrachtungsweise gelten eine ebene Fläche und eine auf eine Ebene abwickelbare Fläche, etwa eine Zylinderfläche, eine Kegelfläche und dergleichen, als ihrem Wesen nach identisch. Die angemessene Methode, die Beschaffenheit einer so betrachteten Fläche allgemein auszudrücken, beruht stets auf der Formel \(\BareRoot{(E\mathrm{d}p^2+2F\mathrm{d}p\mathord{.}\mathrm{d}q+G\mathrm{d}q^2)}\), die das Element mit den beiden Unbestimmten \(p,q\) verknüpft. Bevor wir diesen Gegenstand weiterverfolgen, müssen wir jedoch die Grundsätze der Theorie der kürzesten Linien auf einer gegebenen gekrümmten Fläche voranstellen.
\Article{14}
Eine gekrümmte Linie im Raum wird im Allgemeinen dadurch angegeben, dass die Koordinaten \(x,y,z\) ihrer einzelnen Punkte als Funktionen einer einzigen Veränderlichen dargestellt werden, die wir mit \(w\) bezeichnen. Die Länge einer solchen Linie von einem beliebigen Anfangspunkt bis zu dem Punkt mit den Koordinaten \(x,y,z\) wird durch das Integral ausgedrückt
\[
\int\mathrm{d}w\mathord{.}\BareRoot{((\tfrac{\mathrm{d}x}{\mathrm{d}w})^2+(\tfrac{\mathrm{d}y}{\mathrm{d}w})^2+(\tfrac{\mathrm{d}z}{\mathrm{d}w})^2)}
\]
Nehmen wir an, dass die Lage der gekrümmten Linie eine unendlich kleine Änderung erfährt, sodass die Koordinaten der einzelnen Punkte die Variationen \(\delta x,\delta y,\delta z\) erhalten, so ergibt sich für die Variation der gesamten Länge
\[
=\int\frac{\mathrm{d}x\mathord{.}\mathrm{d}\delta x+\mathrm{d}y\mathord{.}\mathrm{d}\delta y+\mathrm{d}z\mathord{.}\mathrm{d}\delta z}{\BareRoot{(\mathrm{d}x^2+\mathrm{d}y^2+\mathrm{d}z^2)}}
\]
Diesen Ausdruck bringen wir in die folgende Form:
\[
\begin{aligned}
&\frac{\mathrm{d}x\mathord{.}\delta x+\mathrm{d}y\mathord{.}\delta y+\mathrm{d}z\mathord{.}\delta z}{\BareRoot{(\mathrm{d}x^2+\mathrm{d}y^2+\mathrm{d}z^2)}}\\
&-\int(\delta x\mathord{.}\mathrm{d}\frac{\mathrm{d}x}{\BareRoot{(\mathrm{d}x^2+\mathrm{d}y^2+\mathrm{d}z^2)}}+\delta y\mathord{.}\mathrm{d}\frac{\mathrm{d}y}{\BareRoot{(\mathrm{d}x^2+\mathrm{d}y^2+\mathrm{d}z^2)}}\\
&\quad+\delta z\mathord{.}\mathrm{d}\frac{\mathrm{d}z}{\BareRoot{(\mathrm{d}x^2+\mathrm{d}y^2+\mathrm{d}z^2)}})
\end{aligned}
\]
Ist die Linie zwischen ihren Endpunkten die kürzeste, so muss offenbar das, was hier unter dem Integralzeichen steht, verschwinden. Soweit die Linie auf einer
